\begin{spacing}{1}
    \noindent The Regime Próprio de Previdência Social (RPPS) — the public sector pension  system in Brazil — is responsible for providing post-employment benefits to  permanent civil servants. In this context, Ordinance MTP No.\ 1,467/2022  establishes parameters and guidelines for these organizations. Its Article 36  sets a minimum criterion for the selection of mortality tables, requiring that  the adopted table present, at the mean age of the insured population, a life  expectancy equal to or greater than that of the reference table (the annual  IBGE table, extrapolated by sex). This study investigates whether the mortality  table selection criterion established by Article 36 is sufficient to ensure that  the mathematical provision calculated with the adopted table is equal to or  greater than that obtained with the reference table. To this end, simulations  were conducted for four distinct age profiles --- young, balanced, old, and  bimodal --- each comprising 100,000 insured individuals of each gender, under a simplified  scenario with a single old-age retirement benefit, uniform initial salary, and  no decrements due to disability or withdrawal. Mathematical provisions were  calculated using the INE, CUP, and PNI actuarial cost methods, applying the  table selection criterion in accordance with the regulatory framework of  Article 36. The results revealed inconsistencies across all age profiles, with  higher incidence in younger populations, particularly for males. The  investigation showed that the asymmetry between survival levels in the  contributory age range (18--64 years) and in the benefit-receipt age range  (65 years and older) is the underlying mechanism. When the comparative table  projects higher survival at contributory ages than at retirement ages, the  increase in the present value of contributions exceeds the increase in the  present value of benefits, yielding a mathematical provision lower than that  of the reference table. To support prior diagnosis, the index $I_F$ is  proposed, which quantifies the differential in survival rates weighted by the  age structure of the insured population. It is concluded that the criterion of  Article 36, by relying exclusively on life expectancy at the mean age as a  comparative parameter, may be insufficient to guarantee that the adopted  mortality table produces a mathematical provision equal to or greater than that  of the reference table across all possible demographic configurations.
\end{spacing}
\vspace{1.0cm}

\noindent \textbf{Keywords}: Mortality tables. Mathematical provisions. Ministerial Ordinance MTP No. 1,467/2022. Actuarial valuation.